% SYNC-ID: introduction
\section{Introduction}\label{sec:introduction}
% CONTENT_DEFERRED_TO_WRITING_PHASE; visible text is non-substantive scaffolding.
\textit{(THIS IS PLACEHOLDER TEXT, PLZ FULFILL IT WHEN WRITING.)}
% SYNC-ID: introduction-placeholder-1
The opening paragraph will eventually establish the approved context and scope. This fixture uses ordinary prose so that the final document can be inspected as a paper rather than as a collection of macros.
% SYNC-ID: introduction-placeholder-2
The next paragraph is reserved for a concise motivation and a source pointer selected during the writing phase. No field-specific statement is intended here.
% SYNC-ID: introduction-placeholder-3
The closing paragraph of this placeholder section will state the approved scope boundary and remain aligned with the Chinese companion.
% SYNC-ID: introduction-placeholder-4
The surrounding paragraphs also reserve the transition from context to a focused paper outline. The transition should be concise, traceable to the approved plan, and free of claims until that plan has been reviewed.
% SYNC-ID: introduction-placeholder-5
For the purposes of this fixture, the introduction also reserves enough space to test a transition sentence, a short scope qualifier, and a pointer to the later sections. The eventual wording should make clear what is in scope, what is deliberately outside scope, and where each statement can be checked. These are writing tasks rather than findings, so this paragraph carries no empirical or theoretical content.
% SYNC-ID: placeholder-figure
% PLACEHOLDER FIGURE SOURCE -- TODO: replace with an approved, language-neutral figure.
\begin{figure*}[t]
  \centering
  \fbox{\parbox[c][1.05in][c]{0.86\textwidth}{\centering\textit{PLACEHOLDER FIGURE* -- TODO: insert the approved overview diagram here.}}}
  \caption{PLACEHOLDER FIGURE -- TODO: replace with a factual caption after the figure provenance is frozen.}
  \label{fig:placeholder-overview}
\end{figure*}

